\scanpage{159}
$K$ corps de nbs.

$A$ anneau des entiers, alg. finie libre sur $\mathbf Z$.

$\uuline{A}$ schéma en anneaux sur $\mathbf Z$

$\uuline{A}^*$ schéma en groupes.

$\uuline{A}^*(\mathbf Z)=\uuline{\text{unités}}=E$.

$E^0\subset E$ \uline{unités} \uline{absolues} d'indice fini de la forme $2^\alpha$ $(\alpha\unsure{<|\leq}r_1)$ dans $E$, égal à $E$ si $K$ est purement imaginaire.

\deleted{On} Le groupe de Galois de $K$ rendu abélien est alors \inserted{une extension de $\operatorname{Pic}(A)$ par} $\uuline{A}^*(\widehat{\mathbf Z})/E^0$ ou encore par
\[
\uuline{A}^*\bigl(\deleted{\unsure{\mathbf A_{\mathbf Q}}}\mathbf Z\times\mathbf R^*\bigr)/\deleted{\text{\unsure{composante}}}\uuline{A}^*(\mathbf Z)
\]
\uline{divisé par la composante connexe}.
\margin{vérifier ce pt.}

Le schéma en groupes \inserted{connexe} sur $\mathbf Z$ associé à $K$ pour la classification des motifs abéliens sur $K$ semble être alors $G^0=\uuline{A}^*/\mathcal R$, où $\mathcal R$ est le sous-schéma en groupe de $\uuline{A}^*$ \uuline{adhérence} de l'union des sections $E^0$ sur $\mathbf Z$. La fibre gén. est l'adhérence \deleted{schéma} de Zariski de $E^0$ dans $\uuline{K}^*$.
\margin{158}
